% Paper 2, Section 7 — The firing-floor bound: theory with a tested scope limit (evidence E5).
% Numbers verified against ledger entries 471f113 (final shrink-H table, 18/18 cells, 3 seeds
% per cell, floor correction), cef2101 (first refutation + contrapositive retraction),
% 159e341 (why raising M fails structurally).
% CLAIMS DISCIPLINE — forbidden in this section, per the skeleton's claims ledger:
%   "the bound is confirmed"; "the contrapositive earns its keep"; any "190x" floor swing
%   (that figure came from two mislabelled floors and is corrected to 12x).
% The correct floor formula is H_b^{-1}(M*log2 K / H); the earlier log2 M reading is superseded.

\section{The firing-floor bound: theory with a tested scope limit}
\label{sec:bound}

Paper~1 \citep{paper1} derives a firing-floor bound for recurrent networks that carry their
state in spikes: if a binary state of width $H$ with mean activity $\rho$ must distinguish
the memory states a task demands, then
%
\begin{equation}
\rho \;\ge\; H_b^{-1}\!\left(\frac{M \log_2 K}{H}\right),
\label{eq:floor}
\end{equation}
%
where $H_b$ is the binary entropy function. The numerator counts \emph{bits}, not symbols:
copying $M$ symbols drawn from an alphabet of size $K$ requires distinguishing $K^M$ states,
so the demand is $M\log_2 K$ bits.\footnote{An earlier reading of this project used
$\log_2 M$, which would have made every cell below unmeasurable (predicted floors
$\sim 10^{-3}$). All floors quoted here use the corrected form.} The bound is the
theoretical reason to expect a spiking recurrent state to be dense, and it motivated the
whole three-route comparison. This section reports what happened when we tried to make it
bind on a network that still learns its task. It does not bind, in either direction, and we
report that as the negative it is.

\subsection{Raising the memory load cannot test it}
\label{sec:bound-M}

The obvious test is to raise $M$ until the floor is high enough to measure against. It fails
structurally, not for want of compute. On the copy task the \texttt{spikestate} variant is at
exactly chance by $M{=}32$ (accuracy $0.0623 \pm 0.0012$ against a chance level of $0.0625$,
three seeds) and again at $M{=}64$ ($0.0623 \pm 0.0005$). At $H{=}256$ the floor only reaches
$0.500$ at $M{=}64$ --- a load at which the network retains nothing. Equation~\eqref{eq:floor}
prices a state that actually carries the sequence, so a network storing nothing is outside its
scope regardless of what activity it happens to exhibit. \emph{The learnability ceiling of a
spiking recurrent state on copy ($M \approx 16$) sits below the load at which the floor begins
to bind ($M{=}64$).} Consistent with this, measured state activity \emph{falls} with load
($0.496 \pm 0.011$ at $M{=}16$, $0.405 \pm 0.030$ at $M{=}32$) rather than rising as the
bound would require --- but we do not read that as evidence against the bound, because the
premise of retention is already violated. Copy columns must never be cited as bound tests.

\subsection{Shrinking the width is the valid design}
\label{sec:bound-H}

The alternative moves the same ratio $M\log_2 K / H$ by shrinking the denominator. Holding
$M{=}16$ (the load a spiking state can still partly learn, $K{=}16$, so a fixed $64$-bit
demand) and sweeping $H \in \{256, 128, 96, 64\}$ moves the predicted floor
$0.042 \to 0.110 \to 0.174 \to 0.500$, a $12\times$ swing, on a task whose difficulty is
unchanged. Both arms were run at three seeds per cell, $30$ epochs, $80$k copy sequences
(Table~\ref{tab:bound}).

Four criteria were registered in the results ledger before any cell of this row landed:
(i)~a \emph{validity gate} --- the test counts only at widths where the digital reference
still learns copy well above chance; (ii)~\emph{confirmation} --- \texttt{spikestate}'s state
activity sits at or above the predicted floor \emph{and rises as $H$ shrinks}, the trend
being the discriminating observable since the $H{=}256$ point already sits far above its own
floor; (iii)~\emph{refutation} --- activity stays flat or falls across the sweep while the
networks still learn, which would make the observed $\rho \approx 0.25$--$0.5$ an
optimization artifact of the LIF state rather than an information-theoretic floor;
(iv)~a stated \emph{most likely outcome} --- retention may die before the floor binds, in
which case the honest verdict is that the bound is not empirically testable with this
architecture at this scale.

\begin{table}[t]
\centering
\caption{Shrink-$H$ bound row. Copy at $L{=}33$ ($M{=}16$, $K{=}16$, i.e.\ a $64$-bit memory
demand), $30$ epochs, $80$k sequences, three seeds per cell; chance accuracy $0.0625$.
Floor is $H_b^{-1}(64/H)$. $\rho$ is measured \texttt{spikestate} recurrent-state activity;
margin kept is the fraction of the digital reference's above-chance accuracy margin that
\texttt{spikestate} retains; the certificate is $H \cdot H_b(\rho)$ bits, to be compared
against the fixed $64$-bit demand.}
\label{tab:bound}
\small
\begin{tabular}{lrrrrrrr}
\toprule
$H$ & params & floor & digital acc & spikestate acc & $\rho$ & margin kept & certificate \\
\midrule
$256$ & $87{,}889$ & $0.042$ & $0.6302 \pm 0.0106$ & $0.1071 \pm 0.0164$ & $0.4960 \pm 0.0106$ & $0.079$ & $256.0$ \\
$128$ & $28{,}113$ & $0.110$ & $0.5944 \pm 0.0023$ & $0.1522 \pm 0.0049$ & $0.3915 \pm 0.0417$ & $0.169$ & $123.6$ \\
$96$  & $18{,}289$ & $0.174$ & $0.5889 \pm 0.0065$ & $0.1570 \pm 0.0063$ & $0.3909 \pm 0.0336$ & $0.179$ & $92.7$ \\
$64$  & $10{,}513$ & $0.500$ & $0.5600 \pm 0.0095$ & $0.1619 \pm 0.0039$ & $0.2496 \pm 0.0252$ & $0.200$ & $51.9$ \\
\bottomrule
\end{tabular}
\end{table}

\paragraph{The validity gate passes everywhere.}
The digital reference keeps $9$--$10\times$ chance accuracy across an $8.4\times$ parameter
range ($0.630 \to 0.560$ from $H{=}256$ to $H{=}64$). Copy at $M{=}16$ is therefore not
width-limited anywhere in $H \in [64, 256]$: shrinking $H$ moves the predicted floor by
$12\times$ while barely changing what the task demands of a network that can do it. That is
precisely the property raising $M$ could not deliver, and it makes all four cells legitimate
tests rather than capacity-limited ones.

\paragraph{Criterion (ii) is refuted, by anti-correlation rather than by flatness.}
As the predicted floor rises $12\times$, measured state activity \emph{falls} by half:
$0.496 \to 0.392 \to 0.391 \to 0.250$. The $H{=}128$ and $H{=}96$ points are statistically
indistinguishable, and at $H{=}64$ the measured activity lands \emph{below} its own predicted
floor. This is criterion~(iii)'s shape. Measured LIF state activity in this architecture is
\textbf{width-determined, not information-determined}.

\subsection{The anti-capacity finding}
\label{sec:bound-anticapacity}

The sweep also produces an ordering that no information bottleneck can generate.
Inverting~\eqref{eq:floor} turns a measured $(\rho, H)$ into a capacity certificate
$H \cdot H_b(\rho)$: the most bits the spiking state could be carrying. Across the row that
certificate falls $256.0 \to 123.6 \to 92.7 \to 51.9$ bits against a fixed $64$-bit demand,
i.e.\ from a $4.0\times$ surplus to a $0.81\times$ deficit. \emph{Accuracy moves the other
way}: margin kept rises monotonically $0.079 \to 0.169 \to 0.179 \to 0.200$, and paired
$\Delta$bpc against the digital reference improves in lockstep
($+2.509 \pm 0.060 \to +2.079 \pm 0.016 \to +2.050 \pm 0.023 \to +1.829 \pm 0.051$). The most
over-provisioned network is the worst one. A capacity constraint cannot produce that
ordering; an optimization pathology can. Spiking-state failure on precise recall is a
trainability failure, and the bound does not explain it.

\paragraph{A retracted intermediate reading, kept for the record.}
When only the $H{=}64$ cell had landed, its $12$-bit certificate deficit alongside a large
accuracy shortfall looked like the bound earning its keep in contrapositive form --- predicted
incapacity, observed failure. The $H{=}96$ cell killed that reading: the certificate there
shows a $1.4\times$ \emph{surplus} and \texttt{spikestate} nevertheless fails essentially
identically ($0.157$ vs.\ $0.162$ accuracy). The match at $H{=}64$ was coincidental, and the
certificate is necessary-but-not-sufficient and non-diagnostic at this scale. We state this
because the intermediate reading was recorded before it was refuted, and because the
refutation is the more informative result.

\subsection{Verdict}
\label{sec:bound-verdict}

Equation~\eqref{eq:floor} is \textbf{not empirically validated by this architecture at this
scale, in either direction}. We could not construct a regime where the firing-floor bound
binds on a network that still learns the task: raising the memory load kills retention before
the floor engages, and shrinking the width engages the floor but finds measured activity
anti-correlated with it across a $12\times$ swing, with accuracy improving as predicted
capacity shrinks. The bound stays in this paper as theory with a stated and now \emph{tested}
scope limit --- it applies to a spiking recurrent state that genuinely carries the task's
memory, and we were unable to build one.

Two things survive this negative and are worth separating from it. The empirical firing
ceiling of Paper~1 --- a recurrent spiking char-LM that will not sparsify below
$\approx 50\%$ --- is a measurement, not a prediction of~\eqref{eq:floor}, and is untouched.
And the practical design conclusion that a compressed recurrent state should not be carried
in spikes is established here empirically rather than theoretically: Section~\ref{sec:copy}
shows a $1$-bit recurrent state sitting at exactly chance on precise recall at two memory
loads while being the cheapest variant by the energy proxy. The bound was the reason to
expect that; it is not the reason to believe it.
